\section{Viscoelastic two-dimensional bubble}

\status{Not measured. This is the section that separates this paper from a
reproduction of the Newtonian benchmark, and it is also the one with the least
behind it. On the VOF side there are two short exploratory runs --- an
Oldroyd-B and a gPTT configuration, 101 steps each, 20 frames --- which are
enough to show the path executes and not enough to state anything. On the
front-tracking side there are viscoelastic input decks but no production run.
A matched pair has to be run before this section exists.}

\subsection{Case definition}

\missing{the viscoelastic case, chosen so that it is the Newtonian case of
Section~\ref{sec:newtonian} with one parameter changed. Keeping the geometry
and the density and viscosity ratios fixed is what makes the comparison
attributable to rheology; changing the case at the same time as the rheology
would confound the two.}

\subsection{What the Newtonian case cannot show}

\missing{the argument, once there are measurements to carry it. The
three-dimensional companion found that the discriminating feature between the
two interface representations was a \emph{shape} feature --- a trailing cusp
that a cell-wise planar reconstruction cannot resolve below the cell while an
explicit surface keeps resolving below it. Whether the planar analogue exists,
and whether it discriminates, is the question this section has to answer.}

\subsection{Uniform mesh}

\missing{the matched pair: front-tracking and VOF, same deck, same time step,
at two resolutions. Model and parameters to be fixed with the case definition.}

\subsection{Adaptive mesh}

\missing{depends on the uniform pair and on the adaptive criterion.}
